\documentclass{article}
\usepackage[numbers,sort&compress]{natbib}
\usepackage{hyperref}
\title{Incremental Typesetting and its Bibliography}
\author{A. Author}
\begin{document}
\maketitle

\section{Introduction}\label{sec:intro}
The line-breaking algorithm follows \citet{knuth81}, and the page builder
the description in \citet{knuth84}. Earlier systems recomputed everything
after every keystroke \citep{lamport94}; later ones kept what they could
\citep{acar05,hammer14}. Section~\ref{sec:method} says how this one works.

Incremental computation has a long history in compilers and editors, and
the cost model that matters here is the one an author feels: the time from
a keystroke to the page that shows it.

\section{Method}\label{sec:method}
Checkpoints are taken at page boundaries, as in \citet{acar05}; a restart
point is the newest checkpoint before the first changed read, and a run
converges when its state equals the old run's \citep[ch.~3]{hammer14}.
Hyphenation follows \citet{liang83}.

Typesetting a long document again after a one-letter edit takes as long as
the first run; taking it from a checkpoint takes as long as the pages that
changed. The convergence test compares the whole state, so a page after the
edit that is the same as before is not typeset again.

\section{Results}\label{sec:results}
The measurements in Section~\ref{sec:method} agree with the model of
\citet{knuth84} for paragraphs and with that of \citet{plass81} for pages.
See also \citet{lamport94} for the macro layer.

\bibliographystyle{plainnat}
\bibliography{refs}
\end{document}
